\documentclass[10pt,a4paper,sans]{moderncv}

\moderncvstyle{classic}
\moderncvcolor{blue}

\usepackage[scale=0.85]{geometry}
\usepackage[brazil]{babel}
\usepackage[utf8]{inputenc}

\name{Patrick Serrano}{}
\email{patrickcmserrano@gmail.com}
\social[linkedin]{patrickcmserrano}
\social[github]{patrickcmserrano}

\begin{document}
\makecvtitle

\vspace{-1em}
\recipient{Buzzlabs}{leandro@buzzlabs.com.br}
\date{\today}
\opening{Prezado Leandro,}
\closing{Fico à disposição para uma conversa,}
\makelettertitle

Tenho mais de 7 anos de experiência construindo sistemas financeiros em Clojure/ClojureScript. Recentemente, em conversa com o Wallysson Oliveira, soube da abertura dessa vaga e do ambiente técnico de alta qualidade na Buzzlabs, o que me motivou a submeter minha candidatura. Trabalho fullstack por escolha: o mesmo modelo mental funcional que torna o backend auditável e tolerante a falhas --- imutabilidade, Datomic, Pathom/EQL --- aplica-se diretamente ao frontend quando a stack é ClojureScript. Em todos os projetos do meu currículo entreguei as duas pontas, do schema Datomic ao componente na tela.

Meu projeto pessoal mais recente, o Ark Engine, serviu como um laboratório para arquiteturas de alta performance. Na versão Clojure, explorei o estado da arte: XTDB para persistência bitemporal (essencial para auditabilidade de trades), Redis Streams para ingestão e Polylith para isolamento de domínios como o Risk Guard. Desenvolvi também uma versão experimental em Go para validar cenários de latência extrema e otimização de custos de cloud em escala independente. Essa dualidade me permitiu validar que os princípios de snapshots imutáveis e FSM determinístico são independentes de linguagem, mas que a expressividade de Clojure e o modelo do Datomic/XTDB são imbatíveis para a complexidade de regras de negócio financeiras.

A Buzzlabs me interessa pelo modelo de consultoria: trabalhar em domínios diferentes, entregar soluções completas, comunicar decisões técnicas para clientes, operar com autonomia dentro de uma equipe que valoriza qualidade. É o ambiente em que trabalhei nos últimos 7 anos --- cada projeto no meu currículo foi entregue nesse contexto, do backend ao frontend. Sobre Kafka: não tenho experiência em produção, mas trabalhei com Onyx e NATS JetStream cobrindo os mesmos padrões de streaming distribuído, e é uma lacuna que tenho interesse direto em fechar. Estou disponível para uma conversa quando for conveniente para vocês.

\makeletterclosing

\end{document}
